半无限一维液体起初均处在熔点 \(T_m\)，在 \(t=0\) 后令 \(x=0\) 壁面保持 \(T_0<T_m\)，形成 \(0<x<s(t)\) 的固相。固相密度 \(\rho\)、比热 \(c_p\)、导热率 \(k\) 和潜热 \(L\) 常数，固相扩散系数 \(\alpha=k/(\rho c_p)\)，液相始终保持 \(T_m\)，忽略界面动力学过冷。
(1) 从局部热扩散方程和界面潜热守恒推导相似变量形式与精确误差函数温度场。
(2) 求凝固前沿 \(s(t)\) 的系数满足的超越方程，严格证明唯一正根，解释Stefan数与 \(\sqrt t\) 标度。
(3) 在任意时刻验证壁面释放的热量恰等于凝固潜热与固层降温显热之和，避免把全部壁面热流直接当潜热。
(4) 求 \(\mathrm{Ste}\ll1\) 的前沿渐近公式与 \(\mathrm{Ste}\gg1\) 的主导对数标度；讨论初始瞬间热流发散及模型的物理适用范围。
